% Options for packages loaded elsewhere
\PassOptionsToPackage{unicode}{hyperref}
\PassOptionsToPackage{hyphens}{url}
\PassOptionsToPackage{dvipsnames,svgnames,x11names}{xcolor}
%
\documentclass[
  letterpaper,
  DIV=11,
  numbers=noendperiod]{scrartcl}

\usepackage{amsmath,amssymb}
\usepackage{iftex}
\ifPDFTeX
  \usepackage[T1]{fontenc}
  \usepackage[utf8]{inputenc}
  \usepackage{textcomp} % provide euro and other symbols
\else % if luatex or xetex
  \usepackage{unicode-math}
  \defaultfontfeatures{Scale=MatchLowercase}
  \defaultfontfeatures[\rmfamily]{Ligatures=TeX,Scale=1}
\fi
\usepackage{lmodern}
\ifPDFTeX\else  
    % xetex/luatex font selection
\fi
% Use upquote if available, for straight quotes in verbatim environments
\IfFileExists{upquote.sty}{\usepackage{upquote}}{}
\IfFileExists{microtype.sty}{% use microtype if available
  \usepackage[]{microtype}
  \UseMicrotypeSet[protrusion]{basicmath} % disable protrusion for tt fonts
}{}
\makeatletter
\@ifundefined{KOMAClassName}{% if non-KOMA class
  \IfFileExists{parskip.sty}{%
    \usepackage{parskip}
  }{% else
    \setlength{\parindent}{0pt}
    \setlength{\parskip}{6pt plus 2pt minus 1pt}}
}{% if KOMA class
  \KOMAoptions{parskip=half}}
\makeatother
\usepackage{xcolor}
\usepackage[margin=1in]{geometry}
\setlength{\emergencystretch}{3em} % prevent overfull lines
\setcounter{secnumdepth}{5}
% Make \paragraph and \subparagraph free-standing
\makeatletter
\ifx\paragraph\undefined\else
  \let\oldparagraph\paragraph
  \renewcommand{\paragraph}{
    \@ifstar
      \xxxParagraphStar
      \xxxParagraphNoStar
  }
  \newcommand{\xxxParagraphStar}[1]{\oldparagraph*{#1}\mbox{}}
  \newcommand{\xxxParagraphNoStar}[1]{\oldparagraph{#1}\mbox{}}
\fi
\ifx\subparagraph\undefined\else
  \let\oldsubparagraph\subparagraph
  \renewcommand{\subparagraph}{
    \@ifstar
      \xxxSubParagraphStar
      \xxxSubParagraphNoStar
  }
  \newcommand{\xxxSubParagraphStar}[1]{\oldsubparagraph*{#1}\mbox{}}
  \newcommand{\xxxSubParagraphNoStar}[1]{\oldsubparagraph{#1}\mbox{}}
\fi
\makeatother


\providecommand{\tightlist}{%
  \setlength{\itemsep}{0pt}\setlength{\parskip}{0pt}}\usepackage{longtable,booktabs,array}
\usepackage{calc} % for calculating minipage widths
% Correct order of tables after \paragraph or \subparagraph
\usepackage{etoolbox}
\makeatletter
\patchcmd\longtable{\par}{\if@noskipsec\mbox{}\fi\par}{}{}
\makeatother
% Allow footnotes in longtable head/foot
\IfFileExists{footnotehyper.sty}{\usepackage{footnotehyper}}{\usepackage{footnote}}
\makesavenoteenv{longtable}
\usepackage{graphicx}
\makeatletter
\newsavebox\pandoc@box
\newcommand*\pandocbounded[1]{% scales image to fit in text height/width
  \sbox\pandoc@box{#1}%
  \Gscale@div\@tempa{\textheight}{\dimexpr\ht\pandoc@box+\dp\pandoc@box\relax}%
  \Gscale@div\@tempb{\linewidth}{\wd\pandoc@box}%
  \ifdim\@tempb\p@<\@tempa\p@\let\@tempa\@tempb\fi% select the smaller of both
  \ifdim\@tempa\p@<\p@\scalebox{\@tempa}{\usebox\pandoc@box}%
  \else\usebox{\pandoc@box}%
  \fi%
}
% Set default figure placement to htbp
\def\fps@figure{htbp}
\makeatother

\KOMAoption{captions}{tableheading}
\makeatletter
\@ifpackageloaded{tcolorbox}{}{\usepackage[skins,breakable]{tcolorbox}}
\@ifpackageloaded{fontawesome5}{}{\usepackage{fontawesome5}}
\definecolor{quarto-callout-color}{HTML}{909090}
\definecolor{quarto-callout-note-color}{HTML}{0758E5}
\definecolor{quarto-callout-important-color}{HTML}{CC1914}
\definecolor{quarto-callout-warning-color}{HTML}{EB9113}
\definecolor{quarto-callout-tip-color}{HTML}{00A047}
\definecolor{quarto-callout-caution-color}{HTML}{FC5300}
\definecolor{quarto-callout-color-frame}{HTML}{acacac}
\definecolor{quarto-callout-note-color-frame}{HTML}{4582ec}
\definecolor{quarto-callout-important-color-frame}{HTML}{d9534f}
\definecolor{quarto-callout-warning-color-frame}{HTML}{f0ad4e}
\definecolor{quarto-callout-tip-color-frame}{HTML}{02b875}
\definecolor{quarto-callout-caution-color-frame}{HTML}{fd7e14}
\makeatother
\makeatletter
\@ifpackageloaded{caption}{}{\usepackage{caption}}
\AtBeginDocument{%
\ifdefined\contentsname
  \renewcommand*\contentsname{Table of contents}
\else
  \newcommand\contentsname{Table of contents}
\fi
\ifdefined\listfigurename
  \renewcommand*\listfigurename{List of Figures}
\else
  \newcommand\listfigurename{List of Figures}
\fi
\ifdefined\listtablename
  \renewcommand*\listtablename{List of Tables}
\else
  \newcommand\listtablename{List of Tables}
\fi
\ifdefined\figurename
  \renewcommand*\figurename{Figure}
\else
  \newcommand\figurename{Figure}
\fi
\ifdefined\tablename
  \renewcommand*\tablename{Table}
\else
  \newcommand\tablename{Table}
\fi
}
\@ifpackageloaded{float}{}{\usepackage{float}}
\floatstyle{ruled}
\@ifundefined{c@chapter}{\newfloat{codelisting}{h}{lop}}{\newfloat{codelisting}{h}{lop}[chapter]}
\floatname{codelisting}{Listing}
\newcommand*\listoflistings{\listof{codelisting}{List of Listings}}
\makeatother
\makeatletter
\makeatother
\makeatletter
\@ifpackageloaded{caption}{}{\usepackage{caption}}
\@ifpackageloaded{subcaption}{}{\usepackage{subcaption}}
\makeatother

\usepackage{bookmark}

\IfFileExists{xurl.sty}{\usepackage{xurl}}{} % add URL line breaks if available
\urlstyle{same} % disable monospaced font for URLs
\hypersetup{
  pdftitle={Module 1: Schooling for Diverse Populations in American Society},
  pdfauthor={Nathan Alexander, PhD},
  colorlinks=true,
  linkcolor={blue},
  filecolor={Maroon},
  citecolor={Blue},
  urlcolor={Blue},
  pdfcreator={LaTeX via pandoc}}


\title{Module 1: Schooling for Diverse Populations in American Society}
\usepackage{etoolbox}
\makeatletter
\providecommand{\subtitle}[1]{% add subtitle to \maketitle
  \apptocmd{\@title}{\par {\large #1 \par}}{}{}
}
\makeatother
\subtitle{EDUC 219}
\author{Nathan Alexander, PhD}
\date{}

\begin{document}
\maketitle

\renewcommand*\contentsname{Table of contents}
{
\hypersetup{linkcolor=}
\setcounter{tocdepth}{3}
\tableofcontents
}

\section{Module overview}\label{module-overview}

This module introduces EDUC 219 and establishes the key concepts,
historical questions, and shared inquiry practices that will guide the
course. We begin by examining what it means to study schooling for
diverse populations in American society.

The module asks students to connect historical thinking, educational
inquiry, and personal reflection. We will consider how diversity, power,
institutions, and historical memory shape educational experiences and
opportunities.

\subsection{Key concepts and
inquiries}\label{key-concepts-and-inquiries}

\begin{itemize}
\tightlist
\item
  Schooling and education
\item
  Diversity and multiculturalism
\item
  Social justice
\item
  Educational opportunity and unequal access
\item
  Democracy and citizenship
\item
  Historical thinking and historical significance
\item
  Archives, artifacts, and primary sources
\item
  Power, privilege, prejudice, and discrimination
\item
  Professional identity and agency
\end{itemize}

\subsection{Historical artifacts}\label{historical-artifacts}

Historical artifacts and primary sources help us investigate how people,
institutions, communities, policies, and events shaped educational
experiences. In this course, historical inquiry is not only about
identifying facts from the past; it is also about asking how evidence,
interpretation, and historical context shape what we understand about
schools and society.

\begin{tcolorbox}[enhanced jigsaw, colframe=quarto-callout-note-color-frame, rightrule=.15mm, opacitybacktitle=0.6, leftrule=.75mm, colback=white, breakable, coltitle=black, arc=.35mm, bottomrule=.15mm, titlerule=0mm, left=2mm, colbacktitle=quarto-callout-note-color!10!white, title=\textcolor{quarto-callout-note-color}{\faInfo}\hspace{0.5em}{Historical thinking}, toprule=.15mm, bottomtitle=1mm, toptitle=1mm, opacityback=0]

Historical thinking is the practice of using evidence, context,
chronology, and interpretation to understand the past and its
relationships to the present and future.

Historical thinking asks us to consider why we examine the past in order
to better understand the world around us and the world ahead of us. For
educators and informed citizens, it helps us interpret institutions,
policies, social conditions, and unequal opportunities with greater
care.

\end{tcolorbox}

\begin{tcolorbox}[enhanced jigsaw, colframe=quarto-callout-tip-color-frame, rightrule=.15mm, opacitybacktitle=0.6, leftrule=.75mm, colback=white, breakable, coltitle=black, arc=.35mm, bottomrule=.15mm, titlerule=0mm, left=2mm, colbacktitle=quarto-callout-tip-color!10!white, title=\textcolor{quarto-callout-tip-color}{\faLightbulb}\hspace{0.5em}{Primary source}, toprule=.15mm, bottomtitle=1mm, toptitle=1mm, opacityback=0]

A primary source is material created during the historical period being
studied or by someone directly connected to the event, experience,
institution, or idea under investigation.

Examples include court opinions, laws, letters, newspaper articles,
photographs, speeches, school records, meeting minutes, maps, oral
histories, and archival documents.

\end{tcolorbox}

\begin{tcolorbox}[enhanced jigsaw, colframe=quarto-callout-tip-color-frame, rightrule=.15mm, opacitybacktitle=0.6, leftrule=.75mm, colback=white, breakable, coltitle=black, arc=.35mm, bottomrule=.15mm, titlerule=0mm, left=2mm, colbacktitle=quarto-callout-tip-color!10!white, title=\textcolor{quarto-callout-tip-color}{\faLightbulb}\hspace{0.5em}{Historical artifact}, toprule=.15mm, bottomtitle=1mm, toptitle=1mm, opacityback=0]

A historical artifact is an object, document, image, record, or other
source that provides evidence about people, institutions, events, and
social life in the past.

For this course, artifacts may directly concern education or may provide
broader historical context for understanding schools, law, communities,
race, citizenship, and educational opportunity.

\end{tcolorbox}

\begin{tcolorbox}[enhanced jigsaw, colframe=quarto-callout-important-color-frame, rightrule=.15mm, opacitybacktitle=0.6, leftrule=.75mm, colback=white, breakable, coltitle=black, arc=.35mm, bottomrule=.15mm, titlerule=0mm, left=2mm, colbacktitle=quarto-callout-important-color!10!white, title=\textcolor{quarto-callout-important-color}{\faExclamation}\hspace{0.5em}{Major historical event}, toprule=.15mm, bottomtitle=1mm, toptitle=1mm, opacityback=0]

A major historical event is an event that meets a stated set of criteria
for historical significance.

As part of the timeline activity, your group will develop and apply its
own definition. The point is not merely to label an event ``important,''
but to explain what makes it historically significant, to whom, at what
scale, and with what consequences.

\end{tcolorbox}

\subsection{Guiding questions}\label{guiding-questions}

\begin{itemize}
\tightlist
\item
  What is the difference between education and schooling?
\item
  What purposes should schools serve in a democratic society?
\item
  How do race, culture, language, gender, income, and ability shape
  educational experience and opportunity?
\item
  How have historical struggles over diversity and inequality shaped
  American schools?
\item
  What is the relationship among equal access, justice, democracy, and
  teaching?
\item
  How do power, privilege, prejudice, discrimination, and systems of
  oppression affect schools and communities?
\item
  Why does historical thinking matter for educators, students, and
  informed citizens?
\item
  What kinds of sources can help us understand educational history?
\item
  What kind of professional identity and agency do I want to develop?
\end{itemize}

\section{Notes}\label{notes}

\subsection{Schooling and education}\label{schooling-and-education}

Although people often use the terms interchangeably, \textbf{education}
and \textbf{schooling} are not identical.

\begin{tcolorbox}[enhanced jigsaw, colframe=quarto-callout-note-color-frame, rightrule=.15mm, opacitybacktitle=0.6, leftrule=.75mm, colback=white, breakable, coltitle=black, arc=.35mm, bottomrule=.15mm, titlerule=0mm, left=2mm, colbacktitle=quarto-callout-note-color!10!white, title=\textcolor{quarto-callout-note-color}{\faInfo}\hspace{0.5em}{Education}, toprule=.15mm, bottomtitle=1mm, toptitle=1mm, opacityback=0]

Education is the broad process through which people develop knowledge,
skills, values, identities, relationships, and ways of understanding the
world. Education occurs in schools, families, communities, workplaces,
religious institutions, cultural organizations, social movements, and
everyday life.

\end{tcolorbox}

\begin{tcolorbox}[enhanced jigsaw, colframe=quarto-callout-note-color-frame, rightrule=.15mm, opacitybacktitle=0.6, leftrule=.75mm, colback=white, breakable, coltitle=black, arc=.35mm, bottomrule=.15mm, titlerule=0mm, left=2mm, colbacktitle=quarto-callout-note-color!10!white, title=\textcolor{quarto-callout-note-color}{\faInfo}\hspace{0.5em}{Schooling}, toprule=.15mm, bottomtitle=1mm, toptitle=1mm, opacityback=0]

Schooling refers more specifically to formal, organized educational
experiences provided through institutions such as public schools,
private schools, colleges, universities, and other credentialing
systems.

Schooling includes curriculum, assessment, policies, schedules, rules,
professional roles, funding structures, and institutional decisions that
shape access to learning.

\end{tcolorbox}

\subsubsection{Questions for discussion}\label{questions-for-discussion}

\begin{itemize}
\tightlist
\item
  Can a person be educated outside of school?
\item
  What do schools provide that other educational settings may not
  provide?
\item
  What kinds of knowledge, histories, identities, and ways of being are
  recognized or ignored in formal schooling?
\item
  When do schooling and education support one another? When might they
  conflict?
\end{itemize}

\subsection{Diversity}\label{diversity}

Diversity describes differences among people and communities. In EDUC
219, we will consider how multiple forms of diversity shape educational
experience, opportunity, representation, and institutional decision
making.

\begin{tcolorbox}[enhanced jigsaw, colframe=quarto-callout-tip-color-frame, rightrule=.15mm, opacitybacktitle=0.6, leftrule=.75mm, colback=white, breakable, coltitle=black, arc=.35mm, bottomrule=.15mm, titlerule=0mm, left=2mm, colbacktitle=quarto-callout-tip-color!10!white, title=\textcolor{quarto-callout-tip-color}{\faLightbulb}\hspace{0.5em}{Diversity}, toprule=.15mm, bottomtitle=1mm, toptitle=1mm, opacityback=0]

Diversity includes differences related to race, ethnicity, culture,
language, gender, sexuality, income, religion, nationality, immigration
history, disability, ability, family structure, geography, and other
social locations.

Diversity is not simply a list of individual characteristics. It also
concerns how institutions respond to differences and how social systems
distribute recognition, resources, power, and opportunity.

\end{tcolorbox}

\subsubsection{Diversity and educational
opportunity}\label{diversity-and-educational-opportunity}

Race, culture, language, gender, income, and ability can shape:

\begin{itemize}
\tightlist
\item
  Access to schools, courses, teachers, learning materials, and
  technology.
\item
  How students and families are perceived, categorized, or treated.
\item
  Curriculum content and whose histories, languages, knowledge, and
  experiences are represented.
\item
  Discipline, special education, tracking, testing, and other
  institutional practices.
\item
  Students' experiences of belonging, safety, dignity, participation,
  and academic opportunity.
\end{itemize}

\subsection{Multiculturalism, diversity, and social
justice}\label{multiculturalism-diversity-and-social-justice}

These terms are related, but they do not mean the same thing.

\begin{tcolorbox}[enhanced jigsaw, colframe=quarto-callout-note-color-frame, rightrule=.15mm, opacitybacktitle=0.6, leftrule=.75mm, colback=white, breakable, coltitle=black, arc=.35mm, bottomrule=.15mm, titlerule=0mm, left=2mm, colbacktitle=quarto-callout-note-color!10!white, title=\textcolor{quarto-callout-note-color}{\faInfo}\hspace{0.5em}{Multiculturalism}, toprule=.15mm, bottomtitle=1mm, toptitle=1mm, opacityback=0]

Multiculturalism is an approach that recognizes, studies, and values the
histories, cultures, perspectives, and contributions of multiple groups
in society.

In education, multicultural approaches may shape curriculum, texts,
classroom discussion, school programming, and institutional practices.

\end{tcolorbox}

\begin{tcolorbox}[enhanced jigsaw, colframe=quarto-callout-note-color-frame, rightrule=.15mm, opacitybacktitle=0.6, leftrule=.75mm, colback=white, breakable, coltitle=black, arc=.35mm, bottomrule=.15mm, titlerule=0mm, left=2mm, colbacktitle=quarto-callout-note-color!10!white, title=\textcolor{quarto-callout-note-color}{\faInfo}\hspace{0.5em}{Social justice}, toprule=.15mm, bottomtitle=1mm, toptitle=1mm, opacityback=0]

Social justice concerns the fair distribution of rights, resources,
recognition, opportunity, and participation. In education, social
justice asks how schools can identify and challenge inequities rather
than simply treating unequal conditions as normal or inevitable.

\end{tcolorbox}

\begin{tcolorbox}[enhanced jigsaw, colframe=quarto-callout-important-color-frame, rightrule=.15mm, opacitybacktitle=0.6, leftrule=.75mm, colback=white, breakable, coltitle=black, arc=.35mm, bottomrule=.15mm, titlerule=0mm, left=2mm, colbacktitle=quarto-callout-important-color!10!white, title=\textcolor{quarto-callout-important-color}{\faExclamation}\hspace{0.5em}{Equity and equality}, toprule=.15mm, bottomtitle=1mm, toptitle=1mm, opacityback=0]

\textbf{Equality} generally means providing the same resources, rules,
or treatment to everyone.

\textbf{Equity} means responding to unequal conditions, barriers, and
needs in ways intended to make genuine access, participation, and
opportunity possible.

The distinction matters because identical treatment does not always
produce fair or equitable outcomes.

\end{tcolorbox}

\subsection{Power, privilege, prejudice, and
discrimination}\label{power-privilege-prejudice-and-discrimination}

Understanding schooling in American society requires attention to how
social power operates. Schools do not exist apart from their political,
economic, legal, cultural, and historical contexts.

\begin{tcolorbox}[enhanced jigsaw, colframe=quarto-callout-note-color-frame, rightrule=.15mm, opacitybacktitle=0.6, leftrule=.75mm, colback=white, breakable, coltitle=black, arc=.35mm, bottomrule=.15mm, titlerule=0mm, left=2mm, colbacktitle=quarto-callout-note-color!10!white, title=\textcolor{quarto-callout-note-color}{\faInfo}\hspace{0.5em}{Power}, toprule=.15mm, bottomtitle=1mm, toptitle=1mm, opacityback=0]

Power is the ability to shape decisions, define what counts as
knowledge, distribute resources, establish rules, influence
institutions, and affect the conditions of other people's lives.

\end{tcolorbox}

\begin{tcolorbox}[enhanced jigsaw, colframe=quarto-callout-note-color-frame, rightrule=.15mm, opacitybacktitle=0.6, leftrule=.75mm, colback=white, breakable, coltitle=black, arc=.35mm, bottomrule=.15mm, titlerule=0mm, left=2mm, colbacktitle=quarto-callout-note-color!10!white, title=\textcolor{quarto-callout-note-color}{\faInfo}\hspace{0.5em}{Privilege}, toprule=.15mm, bottomtitle=1mm, toptitle=1mm, opacityback=0]

Privilege refers to unearned advantages or reduced barriers that
individuals or groups may experience because of their social position
within broader systems of power.

\end{tcolorbox}

\begin{tcolorbox}[enhanced jigsaw, colframe=quarto-callout-note-color-frame, rightrule=.15mm, opacitybacktitle=0.6, leftrule=.75mm, colback=white, breakable, coltitle=black, arc=.35mm, bottomrule=.15mm, titlerule=0mm, left=2mm, colbacktitle=quarto-callout-note-color!10!white, title=\textcolor{quarto-callout-note-color}{\faInfo}\hspace{0.5em}{Prejudice}, toprule=.15mm, bottomtitle=1mm, toptitle=1mm, opacityback=0]

Prejudice is a preconceived judgment, belief, or attitude about a person
or group that is not based on adequate knowledge or individual
experience.

\end{tcolorbox}

\begin{tcolorbox}[enhanced jigsaw, colframe=quarto-callout-note-color-frame, rightrule=.15mm, opacitybacktitle=0.6, leftrule=.75mm, colback=white, breakable, coltitle=black, arc=.35mm, bottomrule=.15mm, titlerule=0mm, left=2mm, colbacktitle=quarto-callout-note-color!10!white, title=\textcolor{quarto-callout-note-color}{\faInfo}\hspace{0.5em}{Discrimination}, toprule=.15mm, bottomtitle=1mm, toptitle=1mm, opacityback=0]

Discrimination is unequal treatment, exclusion, restriction, or harm
directed toward people because of their membership in a social group or
because they are perceived to belong to that group.

\end{tcolorbox}

\begin{tcolorbox}[enhanced jigsaw, colframe=quarto-callout-important-color-frame, rightrule=.15mm, opacitybacktitle=0.6, leftrule=.75mm, colback=white, breakable, coltitle=black, arc=.35mm, bottomrule=.15mm, titlerule=0mm, left=2mm, colbacktitle=quarto-callout-important-color!10!white, title=\textcolor{quarto-callout-important-color}{\faExclamation}\hspace{0.5em}{Systems of oppression}, toprule=.15mm, bottomtitle=1mm, toptitle=1mm, opacityback=0]

Systems of oppression are interconnected institutional, cultural, legal,
economic, and social arrangements that produce and maintain unequal
power, resources, recognition, and opportunity across social groups.

In this course, we will examine how such systems affect schooling,
teaching, learning, policy, access, and educational outcomes.

\end{tcolorbox}

\subsection{Educational opportunity}\label{educational-opportunity}

Educational opportunity is more than simply being enrolled in a school.
It concerns whether students have meaningful access to high-quality
learning conditions, relationships, resources, curriculum,
participation, and pathways.

\begin{tcolorbox}[enhanced jigsaw, colframe=quarto-callout-tip-color-frame, rightrule=.15mm, opacitybacktitle=0.6, leftrule=.75mm, colback=white, breakable, coltitle=black, arc=.35mm, bottomrule=.15mm, titlerule=0mm, left=2mm, colbacktitle=quarto-callout-tip-color!10!white, title=\textcolor{quarto-callout-tip-color}{\faLightbulb}\hspace{0.5em}{Equal access}, toprule=.15mm, bottomtitle=1mm, toptitle=1mm, opacityback=0]

Equal access refers to the principle that students should not be
excluded from educational opportunities because of race, culture,
language, gender, income, ability, disability, or other social
characteristics.

Questions about equal access include who can enter schools, courses,
programs, activities, and decision-making spaces, as well as what
supports are necessary for meaningful participation.

\end{tcolorbox}

\subsubsection{Questions for
discussion}\label{questions-for-discussion-1}

\begin{itemize}
\tightlist
\item
  What should count as an educational opportunity?
\item
  Can separate institutions be equal in practice?
\item
  How do funding, transportation, housing, curriculum, discipline,
  testing, and teacher assignment shape access?
\item
  What responsibilities do educators, policymakers, courts, families,
  and communities have for creating equitable educational conditions?
\end{itemize}

\section{Historical timelines and archival
inquiry}\label{historical-timelines-and-archival-inquiry}

\subsection{From physical to digital
timelines}\label{from-physical-to-digital-timelines}

In this course, you will work with others to create physical timelines
and then contribute to a shared class digital timeline/database.

Your group will:

\begin{itemize}
\tightlist
\item
  Identify and define criteria for a major historical event.
\item
  Select two bookend events within a broader historical period.
\item
  Identify and document two primary-source artifacts situated between
  the bookend events.
\item
  Determine whether entries are classified as major events and whether
  they are education related.
\item
  Connect records, when relevant, to \emph{Plessy v. Ferguson},
  \emph{Brown v. Board of Education}, both cases, or broader historical
  context.
\end{itemize}

\begin{tcolorbox}[enhanced jigsaw, colframe=quarto-callout-note-color-frame, rightrule=.15mm, opacitybacktitle=0.6, leftrule=.75mm, colback=white, breakable, coltitle=black, arc=.35mm, bottomrule=.15mm, titlerule=0mm, left=2mm, colbacktitle=quarto-callout-note-color!10!white, title=\textcolor{quarto-callout-note-color}{\faInfo}\hspace{0.5em}{Documentation guide}, toprule=.15mm, bottomtitle=1mm, toptitle=1mm, opacityback=0]

Use the \href{../documentation.qmd}{historical documentation guide} when
preparing entries for the class timeline and archival-research database.

\end{tcolorbox}

\subsection{Reflection}\label{reflection}

As you begin the timeline activity, consider:

\begin{itemize}
\tightlist
\item
  Why might two groups select different events as historically
  significant?
\item
  What counts as evidence for a claim about historical importance?
\item
  How can a local event be important even if it is not nationally known?
\item
  How might a source be directly related to education without being
  classified as a major historical event?
\item
  How can historical artifacts help us understand the conditions
  surrounding \emph{Plessy}, \emph{Brown}, and the development of
  American schooling?
\end{itemize}

\section{Philosophies of education}\label{philosophies-of-education}

\subsection{Starting questions}\label{starting-questions}

A philosophy of education is not simply a list of preferred classroom
activities or a philosophy of teaching. It is an argument about the
purposes, values, obligations, and possibilities of education.

As you begin drafting your philosophy of education, consider:

\begin{itemize}
\tightlist
\item
  What purposes should schools serve in a democratic society?
\item
  How have historical struggles over racial, cultural, linguistic,
  gender, income, and ability diversity shaped my understanding of
  education?
\item
  What is the relationship among equal access, justice, democracy, and
  teaching?
\item
  How should educators respond to unequal histories and present
  educational conditions?
\item
  What kind of professional identity and agency do I want to develop?
\end{itemize}

\begin{tcolorbox}[enhanced jigsaw, colframe=quarto-callout-important-color-frame, rightrule=.15mm, opacitybacktitle=0.6, leftrule=.75mm, colback=white, breakable, coltitle=black, arc=.35mm, bottomrule=.15mm, titlerule=0mm, left=2mm, colbacktitle=quarto-callout-important-color!10!white, title=\textcolor{quarto-callout-important-color}{\faExclamation}\hspace{0.5em}{Philosophy of education}, toprule=.15mm, bottomtitle=1mm, toptitle=1mm, opacityback=0]

Your final individual philosophy of education should communicate your
developing understanding of the purposes of education and the
relationships among schools, democracy, access, diversity, power, and
social justice.

It is not a philosophy of teaching. Your work should move beyond
preferred instructional strategies to address the broader purposes and
responsibilities of education.

\end{tcolorbox}

\section{Preparation for the trials}\label{preparation-for-the-trials}

\subsection{Looking ahead}\label{looking-ahead}

Later in the semester, students will work in trial groups focused on
\emph{Plessy v. Ferguson} or \emph{Brown v. Board of Education}.

The historical thinking, archival research, documentation, and
philosophical questions introduced in this module will support that
work. Trial groups will use historical evidence to build context,
analyze arguments, represent competing perspectives, and consider the
educational consequences of Supreme Court decisions.

\begin{tcolorbox}[enhanced jigsaw, colframe=quarto-callout-tip-color-frame, rightrule=.15mm, opacitybacktitle=0.6, leftrule=.75mm, colback=white, breakable, coltitle=black, arc=.35mm, bottomrule=.15mm, titlerule=0mm, left=2mm, colbacktitle=quarto-callout-tip-color!10!white, title=\textcolor{quarto-callout-tip-color}{\faLightbulb}\hspace{0.5em}{Start building connections}, toprule=.15mm, bottomtitle=1mm, toptitle=1mm, opacityback=0]

As you encounter an event, artifact, concept, or course reading, ask:

\begin{itemize}
\tightlist
\item
  Does this help explain the historical context of \emph{Plessy} or
  \emph{Brown}?
\item
  Does it illuminate questions about segregation, citizenship, public
  institutions, schooling, or equal access?
\item
  Could it become a source, claim, question, or piece of evidence for a
  trial group?
\end{itemize}

\end{tcolorbox}




\end{document}
